% Actividad 4.1. Teach-back: ingeniería de prompts. Diapositivas.
% Un solo fuente para las dos salidas (ver Makefile):
%   diapositivas.pdf            solo las diapositivas
%   diapositivas-con-notas.pdf  cada diapositiva con sus notas debajo
\ifdefined\connotas
  \documentclass[aspectratio=169,10pt]{beamer}
  \usepackage{pgfpages}
  \setbeameroption{show notes on second screen=bottom}
\else
  \documentclass[aspectratio=169,10pt]{beamer}
\fi

\usepackage[spanish,es-noshorthands]{babel}
\usepackage{fontspec}
% Nimbus Sans es una Helvetica libre, igual que en el documento de contenido.
\setsansfont{Nimbus Sans}
\setmainfont{Nimbus Sans}
% Símbolo del colón, que Nimbus Sans no trae.
\newfontfamily\fuentecolon{Liberation Sans}
\newcommand{\colones}{{\fuentecolon ₡}}
\usepackage{csquotes}
\usepackage{tcolorbox}
\usepackage{tabularx}
\usepackage{array}

% Colores del MIT y uno por cada pieza del prompt (los mismos del PDF).
\definecolor{mitRojo}{HTML}{A31F34}
\definecolor{mitGris}{HTML}{8A8B8C}
\definecolor{fondo}{HTML}{F3F3F3}
\definecolor{texto}{HTML}{1A1A1A}
\definecolor{pRol}{HTML}{A31F34}
\definecolor{pContexto}{HTML}{1F5C99}
\definecolor{pTarea}{HTML}{2E7D32}
\definecolor{pReglas}{HTML}{B35C00}
\definecolor{pEjemplo}{HTML}{6A4C93}
\definecolor{pFormato}{HTML}{5A5B5C}

% === Tema propio, sobrio =====================================================
\setbeamertemplate{navigation symbols}{}
\setbeamercolor{normal text}{fg=texto,bg=white}
\setbeamercolor{structure}{fg=mitRojo}
\setbeamercolor{frametitle}{fg=mitRojo}
\setbeamerfont{frametitle}{size=\LARGE,series=\bfseries}
\setbeamertemplate{itemize item}{\textcolor{mitRojo}{\small$\blacksquare$}}
\setbeamertemplate{enumerate item}{\textcolor{mitRojo}{\bfseries\insertenumlabel.}}
\setlength{\parskip}{0.35em}

% Barra roja arriba, a todo lo ancho.
\setbeamertemplate{headline}{%
  \begin{beamercolorbox}[wd=\paperwidth,ht=4pt,dp=0pt]{}%
    \textcolor{mitRojo}{\rule{\paperwidth}{4pt}}%
  \end{beamercolorbox}}

\setbeamertemplate{frametitle}{%
  \vspace*{0.45cm}%
  \usebeamerfont{frametitle}\usebeamercolor[fg]{frametitle}%
  \insertframetitle\par\vspace*{0.1cm}}

% Pie discreto, en gris.
\setbeamertemplate{footline}{%
  \hspace*{0.9cm}%
  \begin{minipage}[b]{\dimexpr\paperwidth-1.8cm\relax}
    {\color{mitGris}\rule{\linewidth}{0.4pt}}\par\vspace{2pt}
    \tiny\color{mitGris}Teach-back. Ingeniería de prompts\hfill
    Módulo 4. Actividad 4.1\quad\textbar\quad\insertframenumber
  \end{minipage}\hspace*{0.9cm}\par\vspace*{0.3cm}}

\setbeamersize{text margin left=0.9cm,text margin right=0.9cm}

% Página de notas: solo el texto, grande y limpio, para leer al exponer.
\setbeamertemplate{note page}{%
  \vspace*{0.35cm}%
  \begin{minipage}{\dimexpr\paperwidth-2.2cm\relax}
  {\scriptsize\bfseries\color{mitRojo}Notas para la exposición.
   Diapositiva \insertframenumber\par}
  \vspace{0.12cm}{\color{mitGris}\rule{\linewidth}{0.4pt}}\par\vspace{0.1cm}
  \footnotesize\color{texto}\setlength{\parskip}{0.5em}\raggedright
  \insertnote
  \end{minipage}}

% Etiqueta de color con el nombre de una pieza del prompt.
%   \etiqueta{color}{texto}
\newcommand{\etiqueta}[2]{%
  \tcbox[on line,colback=#1,colframe=#1,boxrule=0pt,arc=1.2pt,
    left=3pt,right=3pt,top=1pt,bottom=1pt]{\color{white}\scriptsize\bfseries #2}}

% Tarjeta de una pieza (diapositiva 3).
%   \tarjeta{color}{nombre}{qué es}{en Subaxa}
\newcommand{\tarjeta}[4]{%
  \begin{tcolorbox}[colback=#1!6,colframe=#1,boxrule=0pt,toprule=3pt,
    arc=0pt,left=7pt,right=7pt,top=4pt,bottom=4pt,height=2.3cm,
    width=\linewidth,before upper=\raggedright]
  {\small\bfseries\color{#1}\MakeUppercase{#2}\par}\vspace{3pt}
  {\small #3\par}\vspace{4pt}
  {\scriptsize\color{mitGris}En Subaxa: }{\scriptsize #4\par}
  \end{tcolorbox}}

% Pieza compacta del prompt (diapositiva 5).
%   \pieza{color}{nombre}{texto}
\newcommand{\pieza}[3]{%
  \begin{tcolorbox}[colback=#1!6,colframe=#1,boxrule=0pt,leftrule=2.5pt,
    arc=0pt,left=4pt,right=4pt,top=1pt,bottom=1pt,
    before skip=0pt,after skip=1pt,before upper=\raggedright]
  \fontsize{6.6}{8}\selectfont
  {\tiny\bfseries\textcolor{#1}{\MakeUppercase{#2}}}\enspace #3
  \end{tcolorbox}}

% Mensaje que saldría al cliente (burbuja gris).
%   \mensaje{título}{texto}
\newcommand{\mensaje}[2]{%
  \begin{tcolorbox}[colback=fondo,colframe=fondo,boxrule=0pt,arc=4pt,
    left=6pt,right=6pt,top=3pt,bottom=4pt,before skip=3pt,after skip=0pt,
    before upper=\raggedright]
  {\tiny\bfseries\color{mitGris}#1\par}\vspace{1pt}
  {\fontsize{6.8}{8.4}\selectfont #2\par}
  \end{tcolorbox}}

% Recuadro con borde (columna, idea o pregunta).
\newtcolorbox{caja}[1][mitGris]{colback=white,colframe=#1,boxrule=0.6pt,
  arc=1.5pt,left=7pt,right=7pt,top=5pt,bottom=5pt,before upper=\raggedright}
\newtcolorbox{destacado}{colback=mitRojo!5,colframe=mitRojo,boxrule=0pt,
  leftrule=3pt,arc=0pt,left=8pt,right=8pt,top=5pt,bottom=5pt}

\begin{document}

% === 1. Portada ==============================================================
{
\setbeamertemplate{footline}{}
\begin{frame}
\vspace*{0.1cm}
\begin{minipage}{\textwidth}
{\footnotesize\bfseries\color{mitGris}MIT Professional Education,
Chief Technology Officer Program\par}
\vspace{1.0cm}
{\small Inteligencia Artificial Generativa Aplicada a la Transformación
Digital\par}
\vspace{0.1cm}
{\scriptsize\color{mitGris}Módulo 4. Actividad 4.1. Teach-back\par}
\vspace{0.6cm}
{\fontsize{30}{35}\selectfont\bfseries\color{mitRojo}
Ingeniería de prompts:\\ qué es y por qué importa\par}
\vspace{0.4cm}
{\normalsize Una explicación para el equipo de
customer success managers de Subaxa\par}
\vspace{1.2cm}
{\color{mitGris}\rule{\linewidth}{0.5pt}\par}
\vspace{0.15cm}
{\small\bfseries Carlos Luis Rojas Aragonés\par}
{\scriptsize\color{mitGris}CTO de Subaxa\quad\textbar\quad 23 de noviembre
de 2026\par}
\end{minipage}
\end{frame}
}

\note{%
Hola a todos. Soy Carlos, el CTO de Subaxa. Esto que les voy a mostrar es
mi teach-back para el curso de IA generativa del programa de CTO del MIT
Professional Education. La tarea es explicarle un concepto a alguien de mi
equipo, y los elegí a ustedes.

El tema es la ingeniería de prompts: qué es y por qué importa. No necesitan
saber nada técnico. Voy a explicar cada palabra rara con palabras de todos
los días, y voy a usar un caso que ustedes ven a diario: los mensajes que
escribe nuestro agente de user engagement por correo y WhatsApp.

¿Por qué ustedes? Porque aprueban cada uno de esos mensajes antes de que
salga. Son la última revisión entre la IA y el inversionista. Y al preparar
esto me di cuenta de que son, además, las personas que mejor pueden hacer
que esos mensajes salgan bien desde el inicio.

Son seis diapositivas más. Al final les dejo cinco preguntas que pueden
usar desde mañana cuando aprueben un mensaje, y una plantilla corta para
proponer mejoras.}

% === 2. Qué es un prompt =====================================================
\begin{frame}{Un prompt es el brief que le damos a la IA}
\textbf{Prompt:} las instrucciones por escrito que le damos a la IA para que
haga un trabajo.

\vspace{0.25cm}
\begin{columns}[T,onlytextwidth]
\begin{column}{0.47\textwidth}
\begin{caja}
{\scriptsize\bfseries\color{mitGris}A UN COMPAÑERO NUEVO, BRIEF VAGO\par}
\vspace{4pt}
\enquote{Escribile a un cliente sobre un remate.}

\vspace{6pt}
{\small\color{mitGris}Resultado: algo correcto, pero genérico. No conoce al
cliente ni nuestras reglas.\par}
\end{caja}
\end{column}
\begin{column}{0.47\textwidth}
\begin{caja}[mitRojo]
{\scriptsize\bfseries\color{mitRojo}AL MISMO COMPAÑERO, BRIEF COMPLETO\par}
\vspace{4pt}
\small Quién es el cliente y qué busca. Qué remate salió. Qué no puede
prometer. Un mensaje que ya funcionó. Corto, para WhatsApp.

\vspace{6pt}
{\small\color{mitGris}Resultado: sale bien desde la primera vez.\par}
\end{caja}
\end{column}
\end{columns}

\vspace{0.45cm}
\begin{destacado}
\textbf{La IA es ese compañero nuevo:} escribe rápido y bien, pero solo sabe
lo que le decimos. Si el brief es vago, el resultado es vago.
\end{destacado}
\end{frame}

\note{%
Empiezo por la palabra. Un prompt son las instrucciones por escrito que le
damos a la IA para que haga un trabajo. Cuando digo IA generativa, o
modelo, me refiero al programa que escribe el texto a partir de esas
instrucciones. Y la ingeniería de prompts es el oficio de escribir esas
instrucciones de forma que la IA entregue lo que necesitamos, y no lo que
ella supone.

La forma más fácil de verlo es pensar en un compañero nuevo. Imaginen que
llega alguien al equipo y le piden: \enquote{escribile a un cliente sobre un
remate}. Va a escribir algo correcto, pero genérico. No es que lo haga mal:
es que no conoce al cliente ni nuestras reglas.

Ahora imaginen que a esa misma persona le dicen quién es el cliente, qué
busca y qué remate salió. Le dicen qué no puede prometer. Le muestran un
mensaje que ya funcionó y le piden que sea corto, para WhatsApp. El mensaje
sale bien desde la primera vez.

La IA es ese compañero nuevo. Escribe rápido y bien, pero solo sabe lo que
le decimos. El prompt es el brief. Si el brief es vago, el resultado es
vago. Esta es la primera idea que quiero que se lleven.}

% === 3. Las seis piezas ======================================================
\begin{frame}{Un buen prompt tiene seis piezas}
\relax{\small Cada pieza responde una pregunta que el compañero nuevo nos haría.\par}
\vspace{0.1cm}
\begin{columns}[T,onlytextwidth]
\begin{column}{0.32\textwidth}
\tarjeta{pRol}{Rol}{Quién \enquote{es} la IA en esta tarea.}
  {alguien del equipo que acompaña a inversionistas.}
\vspace{0.05cm}
\tarjeta{pReglas}{Reglas}{Lo que puede y no puede hacer.}
  {no prometer ganancias ni inventar datos.}
\end{column}
\begin{column}{0.32\textwidth}
\tarjeta{pContexto}{Contexto}{Lo que sabemos de la persona.}
  {zonas, tipo de propiedad, presupuesto, qué vio o guardó, etapa.}
\vspace{0.05cm}
\tarjeta{pEjemplo}{Ejemplo}{Un mensaje modelo de cómo suena
  \enquote{bien}.}{un mensaje corto y cercano ya aprobado.}
\end{column}
\begin{column}{0.32\textwidth}
\tarjeta{pTarea}{Tarea}{Qué tiene que hacer, una sola cosa.}
  {avisar de un remate nuevo que calza con su perfil.}
\vspace{0.05cm}
\tarjeta{pFormato}{Formato}{Cómo se entrega.}
  {canal, largo máximo, un solo llamado a la acción.}
\end{column}
\end{columns}
\end{frame}

\note{%
Un buen prompt tiene seis piezas. Cada una responde una pregunta que el
compañero nuevo nos haría el primer día.

El \textbf{rol}: quién \enquote{es} la IA en esta tarea. En nuestro caso,
alguien del equipo de Subaxa que acompaña a inversionistas.

El \textbf{contexto}: lo que sabemos de la persona. Junto con la tarea le
pasamos a la IA la información que necesita: zonas, tipo de propiedad,
presupuesto, qué remates vio o guardó, en qué etapa está como cliente.

La \textbf{tarea}: qué tiene que hacer, una sola cosa. Por ejemplo, avisar
de un remate nuevo que calza con su perfil.

Las \textbf{reglas}: lo que puede y no puede hacer. No prometer ganancias,
no inventar datos del remate, invitar a hablar con el equipo. Pongo
\enquote{no inventar datos} a propósito: la IA a veces inventa un dato que
suena creíble. A eso se le llama alucinación.

El \textbf{ejemplo}: un mensaje modelo de cómo suena \enquote{bien}. Un
mensaje corto, cercano y profesional que ya aprobamos.

Y el \textbf{formato}: cómo se entrega. El canal, el largo máximo y un solo
llamado a la acción.

Fíjense en los colores. Los voy a usar igual en el ejemplo de dentro de dos
diapositivas.}

% === 4. El agente de user engagement ========================================
\begin{frame}{Nuestro caso: el agente de user engagement}
\relax{\small Está en producción. Escribe mensajes personalizados para cada
usuario de la plataforma.\par}
\vspace{0.15cm}
\begin{columns}[T,onlytextwidth]
\begin{column}{0.2\textwidth}
\begin{caja}
{\scriptsize\bfseries\color{mitRojo}CANALES\par}\vspace{4pt}
\small Correo\par WhatsApp
\end{caja}
\end{column}
\begin{column}{0.46\textwidth}
\begin{caja}
{\scriptsize\bfseries\color{mitRojo}DATOS QUE USA\par}\vspace{4pt}
\footnotesize
Preferencias: zonas, tipo de propiedad y presupuesto.\par
Lo que hace en la plataforma: remates que vio, guardó o en los que
participó.\par
Remates nuevos del Boletín Judicial que calzan.\par
Etapa del cliente: nuevo, activo, inactivo.
\end{caja}
\end{column}
\begin{column}{0.28\textwidth}
\begin{caja}
{\scriptsize\bfseries\color{mitRojo}CUÁNDO ESCRIBE\par}\vspace{4pt}
\footnotesize Cuando pasa algo: sale un remate que calza, o el usuario hace o deja
de hacer algo.\par
Y en un resumen periódico.
\end{caja}
\end{column}
\end{columns}

\vspace{0.3cm}
\begin{destacado}
\textbf{Antes de que un mensaje salga, un CSM lo aprueba.} Ustedes son la
última revisión entre la IA y el inversionista.
\end{destacado}
\end{frame}

\note{%
Ahora llevo esto a algo que ustedes conocen. Nuestro agente de user
engagement está en producción. Escribe mensajes por correo y por WhatsApp
para cada usuario de la plataforma.

Para personalizarlos usa cuatro tipos de datos. Las preferencias del
usuario: zonas, tipo de propiedad y presupuesto. Lo que hace en la
plataforma: los remates que vio, que guardó o en los que participó. Los
remates nuevos del Boletín Judicial que calzan con su perfil. Y la etapa
del cliente: nuevo, activo, inactivo.

Escribe en dos momentos. Cuando pasa algo, por ejemplo cuando aparece un
remate que calza, o cuando el usuario hace o deja de hacer algo. Y además
en un resumen periódico personalizado.

Lo importante para hoy está abajo: antes de que un mensaje salga, uno de
ustedes lo aprueba. Todo lo que el agente escribe pasa por sus manos. Por
eso me interesa que entiendan de dónde sale cada mensaje: sale de un
prompt. Y la calidad del mensaje depende, casi siempre, de la calidad de
ese prompt.}

% === 5. Ana: prompt pobre vs. prompt completo ================================
\begin{frame}[t]{\Large El mismo pedido, con dos prompts distintos}
\vspace{-0.25cm}
\begin{columns}[T,onlytextwidth]
\begin{column}{0.31\textwidth}
{\scriptsize\bfseries\color{mitGris}PROMPT POBRE\par}\vspace{2pt}
\begin{tcolorbox}[colback=white,colframe=mitGris,boxrule=0.5pt,arc=1.5pt,
  left=5pt,right=5pt,top=3pt,bottom=3pt,before skip=0pt,after skip=0pt]
\fontsize{7.2}{9}\selectfont Escriba un mensaje de WhatsApp para un usuario
de Subaxa sobre un remate nuevo.
\end{tcolorbox}
\mensaje{Mensaje que produciría (genérico)}{¡Hola! En Subaxa tenemos nuevas
oportunidades de inversión en remates judiciales. ¡No se las pierda! Entre a
la plataforma y descubra las mejores propiedades al mejor precio. ¡Invierta
hoy!}

\vspace{8pt}
\begin{tcolorbox}[colback=mitRojo!5,colframe=mitRojo,boxrule=0pt,
  leftrule=2.5pt,arc=0pt,left=5pt,right=5pt,top=3pt,bottom=3pt,
  before skip=0pt,after skip=0pt,before upper=\raggedright]
\fontsize{6.8}{8.2}\selectfont\textbf{\color{mitRojo}Ejemplo ilustrativo y
simplificado.} No es el prompt real del agente. Ana, el remate y los montos
son ficticios.
\end{tcolorbox}
\end{column}
\begin{column}{0.66\textwidth}
{\scriptsize\bfseries\color{mitRojo}PROMPT CON LAS SEIS PIEZAS\par}\vspace{2pt}
\pieza{pRol}{Rol}{Equipo de Subaxa que acompaña a inversionistas. Cercano,
de \enquote{usted}.}
\pieza{pContexto}{Contexto}{Ana (ficticia), activa. Busca
casas en Heredia, hasta \colones 80\,000\,000. Guardó dos la semana pasada.
Remate: casa de 3 habitaciones en Heredia, base \colones 65\,000\,000,
15 de enero.}
\pieza{pTarea}{Tarea}{Avísele de este remate y por qué podría interesarle.}
\pieza{pReglas}{Reglas}{Solo estos datos. No prometa ganancias ni diga
\enquote{ganga}. Sin direcciones exactas. Invítela a hablar con el equipo.}
\pieza{pEjemplo}{Ejemplo}{\enquote{Hola, Luis. Vimos que le interesan los
lotes en Cartago\ldots{} Si gusta, lo revisamos juntos.}}
\pieza{pFormato}{Formato}{WhatsApp. Máximo 4 líneas. Un llamado a la acción. Sin emojis.}
\mensaje{Mensaje que produciría (personalizado)}{Hola, Ana. Salió en el
Boletín Judicial una casa de 3 habitaciones en Heredia, parecida a las que
guardó la semana pasada. La base es de \colones 65\,000\,000 y el primer
remate es el 15 de enero. Si gusta, el equipo legal y de inversión le
acompaña a revisarla antes de decidir.}
\end{column}
\end{columns}
\end{frame}

\note{%
Este ejemplo es ilustrativo y simplificado. No es el prompt real del
agente. Ana es una usuaria ficticia, y el remate y los montos también son
ficticios. Lo armé para que se vea la diferencia.

A la izquierda está el prompt pobre: \enquote{escriba un mensaje de WhatsApp
para un usuario de Subaxa sobre un remate nuevo}. Es lo que le diríamos al
compañero nuevo con prisa. El mensaje que sale suena a publicidad masiva. No
dice cuál remate, no sabe a quién le habla y además dice \enquote{al mejor
precio}, que es una promesa que no podemos hacer.

A la derecha está el mismo pedido con las seis piezas, con los mismos
colores de antes. El rol: alguien del equipo, cercano y profesional, que
trata de \enquote{usted}. El contexto: Ana busca casas en Heredia, tiene un
presupuesto de hasta 80 millones de colones y la semana pasada guardó dos
casas. Y el remate nuevo: una casa de 3 habitaciones en Heredia con base de
65 millones. La tarea: avisarle de ese remate. Las reglas: solo esos datos,
nada de ganancias ni \enquote{gangas}, sin direcciones exactas. Un ejemplo
de tono. Y el formato: WhatsApp, máximo cuatro líneas, un solo llamado a la
acción.

El mensaje de abajo usa su zona, su tipo de propiedad y lo que hizo en la
plataforma. Está dentro de su presupuesto. No promete nada. Y es corto. La
IA es la misma en los dos casos: lo único que cambió fue el prompt.}

% === 6. Qué cambia para los CSM ==============================================
\begin{frame}{Qué cambia para ustedes}
\relax{\small Cuando un mensaje sale mal, casi siempre falta o falla una pieza.\par}
\vspace{0.3cm}
\begin{columns}[T,onlytextwidth]
\begin{column}{0.47\textwidth}
\begin{caja}
{\scriptsize\bfseries\color{mitGris}CORREGIR EL MENSAJE A MANO\par}
\vspace{5pt}
Arregla \textbf{un} mensaje.\par
\vspace{3pt}
{\small\color{mitGris}El mismo error vuelve a aparecer en el siguiente.\par}
\end{caja}
\end{column}
\begin{column}{0.47\textwidth}
\begin{caja}[mitRojo]
{\scriptsize\bfseries\color{mitRojo}AJUSTAR EL PROMPT UNA VEZ\par}
\vspace{5pt}
Arregla la causa.\par
\vspace{3pt}
{\small\color{mitGris}Mejora \textbf{todos} los mensajes que vienen.\par}
\end{caja}
\end{column}
\end{columns}

\vspace{0.45cm}
\begin{destacado}
\textbf{Ustedes ya tienen el criterio.} Saben cómo se le habla a un cliente
nuevo y cómo a uno inactivo. El prompt es el lugar donde ese criterio queda
por escrito para la IA.
\end{destacado}
\end{frame}

\note{%
¿Qué cambia para ustedes con todo esto? Cuando un mensaje sale mal, casi
siempre falta o falla una pieza. Si falta el contexto, sale genérico. Si
falta una regla, puede prometer algo que no debe. Si falta el formato, sale
largo o en el tono equivocado.

Ante un mensaje malo hay dos caminos. El primero es corregirlo a mano. Eso
arregla un mensaje, pero el mismo error vuelve a aparecer en el siguiente,
y en el otro. El segundo es ajustar el prompt una vez. Eso arregla la
causa, y el ajuste mejora todos los mensajes que vienen.

Cuando el prompt es pobre, ustedes terminan reescribiendo mensajes en lugar
de revisarlos. Cuando el prompt es bueno, aprobar es más rápido, y la
revisión se enfoca en lo que de verdad requiere criterio humano: el tono con
un cliente en particular, o un caso que el prompt no previó.

Y ese criterio ya lo tienen. Ustedes saben cómo se le habla a un cliente
nuevo y cómo a uno inactivo. El prompt es el lugar donde ese criterio queda
por escrito para la IA. Por eso su papel no es solo aprobar: también pueden
mejorar el prompt.}

% === 7. Las cinco preguntas y la plantilla ===================================
\begin{frame}{Qué pueden hacer desde mañana}
\begin{columns}[T,onlytextwidth]
\begin{column}{0.53\textwidth}
{\scriptsize\bfseries\color{mitRojo}AL APROBAR UN MENSAJE, CINCO PREGUNTAS\par}
\vspace{2pt}
\begin{caja}[mitRojo]
\small
\begin{enumerate}\setlength{\itemsep}{3pt}
\item ¿Se nota que es para esta persona?
\item ¿Los datos del remate son correctos y no hay nada inventado?
\item ¿Promete algo que no debe?
\item ¿El tono calza con la etapa del cliente?
\item ¿Tiene el largo del canal y un solo llamado a la acción?
\end{enumerate}
\end{caja}
\end{column}
\begin{column}{0.42\textwidth}
{\scriptsize\bfseries\color{mitGris}SI UNA RESPUESTA ES \enquote{NO}\par}
\vspace{2pt}
\begin{tcolorbox}[colback=fondo,colframe=fondo,arc=1.5pt,left=7pt,
  right=7pt,top=5pt,bottom=5pt]
{\scriptsize\bfseries Plantilla para sugerir un ajuste al prompt\par}
\vspace{4pt}
\scriptsize
\begin{tabularx}{\linewidth}{@{}>{\bfseries}p{1.55cm}>{\raggedright\arraybackslash}X@{}}
Pieza: & \textbf{\color{pReglas}Reglas}\\[3pt]
Qué pasó: & El mensaje dijo \enquote{excelente oportunidad de
  ganancia}.\\[3pt]
Ajuste: & Agregar la regla \enquote{no hablar de ganancias ni de
  rentabilidad esperada}.\\
\end{tabularx}
\end{tcolorbox}
{\scriptsize\color{mitGris}Además de corregir el mensaje.\par}
\end{column}
\end{columns}

\vspace{0.4cm}
\begin{destacado}
\small Un mensaje genérico casi nunca es culpa de la IA: es un prompt que
dejó decisiones sin tomar.
\end{destacado}
\end{frame}

\note{%
Cierro con algo que pueden usar desde mañana. Cuando aprueben un mensaje,
háganse estas cinco preguntas rápidas.

Uno: ¿se nota que es para esta persona? Que mencione su zona, su tipo de
propiedad o lo que hizo en la plataforma. Dos: ¿los datos del remate son
correctos y no hay nada inventado? Tres: ¿promete algo que no debe? Por
ejemplo ganancias, una \enquote{ganga} o certeza del resultado. Cuatro: ¿el
tono calza con la etapa del cliente? No se le habla igual a uno nuevo que a
uno inactivo. Cinco: ¿tiene el largo del canal y un solo llamado a la
acción?

Si una respuesta es \enquote{no}, corrijan el mensaje, como hacen hoy. Pero
además escriban una sugerencia de ajuste con esta plantilla: qué pieza
falló, qué pasó y qué ajuste proponen. En el ejemplo, el mensaje dijo
\enquote{excelente oportunidad de ganancia}; la pieza que falló son las
reglas, y el ajuste es agregar \enquote{no hablar de ganancias ni de
rentabilidad esperada}.

Me quedo con lo que aprendí al preparar esto. Un mensaje genérico casi nunca
es culpa de la IA: es un prompt que dejó decisiones sin tomar. Y las
personas que mejor pueden tomar esas decisiones son ustedes, no yo. Muchas
gracias.}

\end{document}
